\section{Место в области}
\label{sec:related}
\porozhdeno{\Claim{L01}{\istochnik{Bokulich и соавторы сравнивали и настраивали методы таксономического назначения в q2-feature-classifier на модельных сообществах и симуляциях~\cite{bokulich2018}.} Здесь сравниваются готовые модели одного алгоритма на естественных почвах, без оптимизации по эталону.}}

\porozhdeno{\Claim{L02}{\istochnik{GTDB предлагает геномную таксономию бактерий и архей, согласованную с филогенией и нормированную по рангам~\cite{parks2022}.} Наша задача касается назначений коротких ампликонов, а не построения геномной таксономии; эта публикация не удостоверяет обучение использованного файла модели.}
\Claim{L03}{\istochnik{Greengenes2 объединяет геномные и 16S-данные в общем референсном дереве~\cite{mcdonald2024}.} В нашем сравнении используются только ампликонные ASV, без shotgun-данных и построения общего дерева.}}

\porozhdeno{\Claim{L04}{\istochnik{RESCRIPt обеспечивает воспроизводимое получение, обработку и оценку референсов~\cite{robeson2021}.} Здесь базы не конструируются, модели не переобучаются; используются GG2 вместо рассмотренного в RESCRIPt старого Greengenes и буквальные родословные без объединения значимых суффиксов.}}

\porozhdeno{\Claim{L05}{\istochnik{DADA2 выводит точные варианты ампликонных последовательностей с помощью модели ошибок~\cite{callahan2016}.} Мы фиксируем результат денойзинга, не сравнивая его алгоритмы.}
\Claim{L06}{\istochnik{QIIME 2 сохраняет вычислительное происхождение артефактов и использует семантические типы~\cite{bolyen2019}.} Мы применяем эти механизмы для проверки конкретного набора данных, не разрабатывая новую платформу.}}

\porozhdeno{\Claim{L07}{\istochnik{Balvočiūtė и Huson предложили явное сопоставление сущностей разных таксономий~\cite{balvociute2017}.} Мы сравниваем назначения одних ASV и профили образцов, не строя отображения между деревьями и не объединяя неподтверждённые синонимы.}}
